\documentclass[12pt,letterpaper]{article}
\usepackage[utf8]{inputenc}
\usepackage[T1]{fontenc}
\usepackage[spanish]{babel}
\usepackage{amsmath,amssymb}
\usepackage[margin=2.5cm]{geometry}
\usepackage{enumitem}
\setlength{\parindent}{0pt}
\setlength{\parskip}{0.7em}

\begin{document}

\begin{center}
  {\Large\textbf{Resumen: Investigación de Operaciones}}
\end{center}

\textbf{Integrantes del equipo:}\\
Calderón Hernández Yunuel Jatziry\\
Vilchis Fragoso Ximena Rosario\\
Blas Rosas Casandra\\
Moctezuma Juan Mayrin Paola

%--------------------------------------------------------------
\section*{Del problema a un modelo}

La Investigación de Operaciones busca transformar problemas reales de asignación de recursos en modelos matemáticos que permitan analizar y encontrar soluciones. El procedimiento consiste en comprender la situación, identificar las decisiones, construir el modelo y posteriormente interpretar la solución.

Pasos para construir un modelo de programación lineal:
\begin{itemize}
  \item Definir las decisiones y sus unidades.
  \item Asignar una variable de decisión a cada decisión.
  \item Construir la función objetivo, que puede buscar maximizar un beneficio o minimizar un costo.
  \item Establecer las restricciones que representan los recursos, requisitos o condiciones del problema.
  \item Agregar las condiciones de no negatividad.
  \item Revisar que las dimensiones, unidades y significado de cada expresión sean correctos.
\end{itemize}

\subsection*{Función objetivo}

\[ Z = c_1x_1 + c_2x_2 + \cdots + c_nx_n \]

Para maximización:
\[ \text{Max } Z = c_1x_1 + c_2x_2 + \cdots + c_nx_n \]

Para minimización:
\[ \text{Min } Z = c_1x_1 + c_2x_2 + \cdots + c_nx_n \]

Condición de no negatividad:
\[ x_i \geq 0 \]

Un modelo es lineal cuando las variables aparecen solamente con potencia 1 y no se multiplican entre sí.

%--------------------------------------------------------------
\section*{Matrices}

Una matriz organiza información en filas y columnas. Si tiene $m$ filas y $n$ columnas, su dimensión es $m \times n$.
\[ A = (a_{ij})_{m\times n} \]

Para sumar o restar matrices deben tener las mismas dimensiones. Para multiplicarlas, el número de columnas de la primera debe coincidir con el número de filas de la segunda.
\[ A_{m\times n}\, B_{n\times p} = C_{m\times p} \]

La transpuesta se obtiene intercambiando filas por columnas:
\[ A^{T} = (a_{ji}) \]

Un sistema de ecuaciones puede escribirse en forma matricial como:
\[ AX = B \]

Y su matriz aumentada como:
\[ [\,A \mid B\,] \]

El método de Gauss utiliza operaciones elementales sobre las filas para transformar el sistema y facilitar su resolución.

%--------------------------------------------------------------
\section*{Preparar el modelo para Simplex}

Para aplicar el método Simplex, las restricciones deben llevarse a una forma adecuada. Dependiendo del signo se agregan variables de holgura, exceso o, cuando es necesario, variables artificiales.

\subsection*{Restricción menor o igual ($\leq$)}
\[ a^{T}x \leq b \;\rightarrow\; a^{T}x + s = b, \quad s \geq 0 \]
La variable $s$ es una variable de holgura y representa el recurso que no se utiliza.

\subsection*{Restricción mayor o igual ($\geq$)}
\[ a^{T}x \geq b \;\rightarrow\; a^{T}x - e = b, \quad e \geq 0 \]
La variable $e$ es una variable de exceso y representa cuánto se supera un mínimo requerido.

\subsection*{Variable artificial}
\[ a \geq 0 \]
Se utiliza para obtener una base inicial en determinadas restricciones y debe desaparecer de la solución final.

\subsection*{Cálculos de la tabla Simplex}

Cálculo de $Z_j$:
\[ Z_j = \sum_i C_{B_i}\, a_{ij} \]

Costo reducido:
\[ C_j - Z_j \]

En maximización entra la variable con el mayor valor positivo y se termina cuando:
\[ C_j - Z_j \leq 0 \]

En minimización entra la variable con el valor más negativo y se termina cuando:
\[ C_j - Z_j \geq 0 \]

\subsection*{Prueba de la razón}
\[ \theta_i = \frac{b_i}{a_{ik}}, \quad \text{siempre que } a_{ik} > 0 \]
Se elige la menor razón no negativa.

\subsection*{Ejemplo de Gran M}

\[
\begin{aligned}
\text{Min } C &= 6x_1 + 8x_2\\
x_1 + x_2 &= 5000\\
x_1 &\leq 1500\\
x_2 &\geq 750\\
x_1, x_2 &\geq 0
\end{aligned}
\]

Forma estándar:
\[
\begin{aligned}
x_1 + x_2 + a_1 &= 5000\\
x_1 + s_1 &= 1500\\
x_2 - e_2 + a_2 &= 750
\end{aligned}
\]

Función penalizada:
\[ \text{Min } C = 6x_1 + 8x_2 + M a_1 + M a_2 \]

Solución:
\[ x_1 = 1500, \quad x_2 = 3500 \]
\[ C_{\min} = 37000 \]
\[ e_2 = 2750, \quad s_1 = 0 \]

%--------------------------------------------------------------
\section*{Ejemplo completo de maximización}

Función objetivo:
\[ \text{Max } Z = 4x_1 + 6x_2 \]

Sujeto a:
\[
\begin{aligned}
0.5x_1 + x_2 &\leq 600\\
12.5x_1 + 10x_2 &\leq 10000\\
x_1 &\leq 700\\
x_1, x_2 &\geq 0
\end{aligned}
\]

Forma estándar:
\[
\begin{aligned}
0.5x_1 + x_2 + s_1 &= 600\\
12.5x_1 + 10x_2 + s_2 &= 10000\\
x_1 + s_3 &= 700
\end{aligned}
\]

Solución:
\[ x_1 = \frac{1600}{3} \approx 533.33 \]
\[ x_2 = \frac{1000}{3} \approx 333.33 \]
\[ Z_{\max} = \frac{12400}{3} \approx 4133.33 \]

Variables de holgura:
\[ s_1 = 0, \quad s_2 = 0, \quad s_3 = \frac{500}{3} \approx 166.67 \]

\end{document}
